\documentclass[11pt]{article}
\usepackage[utf8]{inputenc}
\usepackage{geometry}
\geometry{a4paper, margin=1in}
\usepackage{hyperref}
\usepackage{listings}
\usepackage{xcolor}
\usepackage{booktabs}
\usepackage{amsmath}

\hypersetup{
    colorlinks=true,
    linkcolor=blue,
    filecolor=magenta,      
    urlcolor=cyan,
    pdftitle={AXI-Lite Register Validation Specification},
    pdfpagemode=FullScreen,
}

\definecolor{codegreen}{rgb}{0,0.6,0}
\definecolor{codegray}{rgb}{0.5,0.5,0.5}
\definecolor{codepurple}{rgb}{0.58,0,0.82}
\definecolor{backcolour}{rgb}{0.95,0.95,0.92}

\lstdefinestyle{mystyle}{
    backgroundcolor=\color{backcolour},   
    commentstyle=\color{codegreen},
    keywordstyle=\color{magenta},
    numberstyle=\tiny\color{codegray},
    stringstyle=\color{codepurple},
    basicstyle=\ttfamily\footnotesize,
    breakatwhitespace=false,         
    breaklines=true,                 
    captionpos=b,                    
    keepspaces=true,                 
    numbers=left,                    
    numbersep=5pt,                  
    showspaces=false,                
    showstringspaces=false,
    showtabs=false,                  
    tabsize=2
}

\lstset{style=mystyle}

\title{AXI-Lite Register Validation Specification\\
\large Lenseup EDGE Artix-7 FPGA Development Board Implementation}
\author{Post-Silicon Verification Team}
\date{June 24, 2026}

\begin{document}

\maketitle

\tableofcontents

\newpage

\section{System Overview}
The AXI-Lite Register Validation system implements an automated framework for testing registers inside the Xilinx Artix-7 FPGA chip. The primary goal is to bridge high-level Python verification scripts on a host computer to physical AXI-Lite registers on the hardware target over a single micro-USB serial link.

\section{Hardware and Port Configuration}
The hardware utilizes the following clocking, reset, and serial properties:
\begin{itemize}
    \item \textbf{FPGA Target}: Xilinx Artix-7 \texttt{xc7a35tftg256-1}.
    \item \textbf{Oscillator}: 50 MHz clock signal routed to Pin \texttt{N11}.
    \item \textbf{Reset Layout}: Active-high push button (\texttt{pb[4]}) on Pin \texttt{M14}. It is logically inverted to active-low (\texttt{rst\_n}) in the top-level Verilog wrapper.
    \item \textbf{Baud Rate}: 115,200 bps.
    \item \textbf{UART Physical Pins}: Receive (RX) connected to Pin \texttt{D4}; Transmit (TX) connected to Pin \texttt{C4}.
\end{itemize}

\section{Serial Command Framing Protocol}
Communication between the host and the FPGA is based on packet transactions:

\subsection{Write Request and Response (Opcode: 'W' / 0x57)}
The host writes a 32-bit word using a 9-byte command frame, and receives a 6-byte response:
\begin{lstlisting}[language=bash, caption=Write Serial Packets]
# Request (9 bytes): [0x57] [Addr31:24] [Addr23:16] [Addr15:8] [Addr7:0] [Data31:24] [Data23:16] [Data15:8] [Data7:0]
# Example: Write 0xDEADBEEF to Address 0x0000000C (SCRATCH Register)
0x57 0x00 0x00 0x00 0x0C 0xDE 0xAD 0xBE 0xEF

# Response (6 bytes): [Status] [Data31:24] [Data23:16] [Data15:8] [Data7:0] [0x57]
0x00 0xDE 0xAD 0xBE 0xEF 0x57
\end{lstlisting}

\subsection{Read Request and Response (Opcode: 'R' / 0x52)}
The host reads a 32-bit word using a 5-byte command frame, and receives a 6-byte response:
\begin{lstlisting}[language=bash, caption=Read Serial Packets]
# Request (5 bytes): [0x52] [Addr31:24] [Addr23:16] [Addr15:8] [Addr7:0]
# Example: Read Address 0x00000100 (DEVICE_ID Register)
0x52 0x00 0x00 0x01 0x00

# Response (6 bytes): [Status] [Data31:24] [Data23:16] [Data15:8] [Data7:0] [0x52]
0x00 0xA5 0xA5 0x00 0x09 0x52
\end{lstlisting}

\newpage

\section{AXI-Lite Register Addressing}
The interconnect routes transactions based on three base blocks. Address bits \texttt{[3:2]} are decoded at each slave interface (address folding), word-aligning all offsets automatically.

\begin{table}[h]
\centering
\caption{System Register Definitions}
\begin{tabular}{@{}llllll@{}}
\toprule
\textbf{Block} & \textbf{Base Address} & \textbf{Register} & \textbf{Offset} & \textbf{Mode} & \textbf{Reset Value} \\ \midrule
\textbf{gpio}  & \texttt{0x00000000}    & \texttt{GPIO\_OUT} & \texttt{0x00}   & RW            & \texttt{0x00000000}  \\
               &                        & \texttt{GPIO\_IN}  & \texttt{0x04}   & RO            & \texttt{0x00000000}  \\
               &                        & \texttt{GPIO\_DIR} & \texttt{0x08}   & RW            & \texttt{0x00000000}  \\
               &                        & \texttt{SCRATCH}   & \texttt{0x0C}   & RW            & \texttt{0x00000000}  \\ \midrule
\textbf{status}& \texttt{0x00000100}    & \texttt{DEVICE\_ID}& \texttt{0x00}   & RO            & \texttt{0xA5A50009}  \\
               &                        & \texttt{STATUS\_FLAGS}& \texttt{0x04}& RO            & \texttt{0x00000000}  \\
               &                        & \texttt{ERROR\_CLEAR}& \texttt{0x08} & WO            & \texttt{0x00000000}  \\
               &                        & \texttt{VERSION}   & \texttt{0x0C}   & RO            & \texttt{0x00010000}  \\ \midrule
\textbf{timer} & \texttt{0x00000200}    & \texttt{TIMER\_CTRL}& \texttt{0x00}  & RW            & \texttt{0x00000000}  \\
               &                        & \texttt{TIMER\_COUNT}& \texttt{0x04} & RO            & \texttt{0x00000000}  \\
               &                        & \texttt{TIMER\_LIMIT}& \texttt{0x08} & RW            & \texttt{0x00000000}  \\
               &                        & \texttt{TIMER\_STATUS}& \texttt{0x0C}& RO            & \texttt{0x00000000}  \\ \bottomrule
\end{tabular}
\end{table}

\section{AXI-Lite Bus Handshaking Channels}
All transfers are completed via \texttt{VALID}/\texttt{READY} handshakes across 5 channels:
\begin{itemize}
    \item \textbf{AW Channel}: \texttt{s\_axi\_awaddr}, \texttt{s\_axi\_awvalid}, and \texttt{s\_axi\_awready}.
    \item \textbf{W Channel}: \texttt{s\_axi\_wdata}, \texttt{s\_axi\_wstrb}, \texttt{s\_axi\_wvalid}, and \texttt{s\_axi\_wready}.
    \item \textbf{B Channel}: \texttt{s\_axi\_bresp}, \texttt{s\_axi\_bvalid}, and \texttt{s\_axi\_bready}.
    \item \textbf{AR Channel}: \texttt{s\_axi\_araddr}, \texttt{s\_axi\_arvalid}, and \texttt{s\_axi\_arready}.
    \item \textbf{R Channel}: \texttt{s\_axi\_rdata}, \texttt{s\_axi\_rresp}, \texttt{s\_axi\_rvalid}, and \texttt{s\_axi\_rready}.
\end{itemize}

\section{Passive Hardware Monitors}
The design incorporates internal passive monitors to aid validation:
\begin{itemize}
    \item \textbf{Protocol Checker}: Latches protocol checker violations on handshaking stability to status flag bit 0.
    \item \textbf{Access Monitor}: Counts read and write bus cycles and logs the address of the last transaction.
    \item \textbf{Reset Recovery Monitor}: Monitors stabilization windows after a hard reset, raising a flag if a transaction is initiated too early.
\end{itemize}

\end{document}
